\section{Related Work}
\label{sec:related}

\paragraph{Hardware Description Languages.}
The limitations of Verilog and VHDL have motivated a rich line of work on higher-level HDLs. Chisel~\citep{bachrach2012chisel} embeds hardware construction in Scala, providing parameterized generators and the FIRRTL intermediate representation that eliminates unintended latches through its lowering passes. Bluespec SystemVerilog~\citep{nikhil2004bluespec} uses guarded atomic actions to prevent race conditions by construction. C$\lambda$ash~\citep{baaij2010clash} leverages Haskell's type system to encode hardware semantics, enabling equational reasoning about circuits. PyMTL~\citep{lockhart2014pymtl} and MyHDL~\citep{decaluwe2004myhdl} prioritize accessibility by embedding hardware in Python. High-level synthesis (HLS) tools~\citep{cong2011hls,canis2011legup} raise the abstraction further by compiling C/C++ to RTL, but sacrifice fine-grained control over micro-architecture. While these approaches improve productivity, none provides the combination of dependent types, formal verification, and synthesizable extraction that Sparkle achieves through its embedding in Lean~4.

\paragraph{Formal Verification of Hardware.}
Formal methods have a long history in hardware verification, from model checking~\citep{clarke2018model} and bounded model checking~\citep{biere2009bounded} to equivalence checking via tools like ABC~\citep{brayton2010abc}. These approaches typically operate on the RTL or gate-level netlist \emph{after} the design is written, treating verification as a separate post-hoc phase. An alternative line of work embeds hardware in proof assistants to enable \emph{verification by construction}. Kami~\citep{choi2017kami} provides a Bluespec-like language in Coq with modular refinement proofs. Silver Oak~\citep{silveroak2021} generates verified hardware from Coq specifications. ACL2 has been used extensively in industrial processor verification~\citep{hunt2017industrial}. Sparkle follows this tradition but targets Lean~4~\citep{moura2021lean4}, whose metaprogramming capabilities enable seamless SystemVerilog extraction alongside formal proofs, and whose dependent type system provides richer compile-time feedback than prior proof-assistant HDLs.

\paragraph{LLMs for Hardware Design.}
The application of LLMs to hardware design has gained significant momentum. Early work by \citet{pearce2020dave} explored generating Verilog from English descriptions. More recently, general-purpose code models such as Codex~\citep{chen2021codex}, Code Llama~\citep{roziere2023codellama}, and StarCoder~2~\citep{li2024starcoder2} have been evaluated on Verilog generation tasks. VeriGen~\citep{thakur2024verigen} fine-tunes LLMs specifically on Verilog corpora, while Chip-Chat~\citep{blocklove2023chip} and ChipGPT~\citep{chang2024chipgpt} explore conversational and multi-step hardware design flows. Benchmarks such as VerilogEval~\citep{liu2023verilogeval} and RTLLM~\citep{lu2024rtllm} provide standardized evaluation suites. A common limitation across this body of work is that LLMs generate Verilog directly, inheriting all of its pitfalls: the generated code may contain width mismatches, unintended latches, or combinational loops that pass syntax checks but fail in synthesis or silicon. Our work addresses this by redirecting the LLM to target a type-safe HDL, where the compiler catches these errors immediately and provides structured feedback for iterative correction.

\paragraph{LLM Agents for Code Generation.}
Tool-augmented LLM agents have shown strong results in software engineering. ReAct~\citep{yao2023react} interleaves reasoning and action in a general agent framework. SWE-agent~\citep{yang2024swe} and SWE-bench~\citep{jimenez2024swebench} demonstrate that LLM agents can resolve real-world GitHub issues by reading code, editing files, and running tests in a feedback loop. Our coding agent follows a similar paradigm---iterative generation, compilation, and error-driven repair---but operates in the hardware domain, where the feedback signal from a dependently-typed compiler is substantially richer than typical software build errors.

\paragraph{LLMs for Theorem Proving.}
LLMs have also been applied to formal theorem proving, with systems like GPT-f~\citep{polu2020generative} and LeanDojo~\citep{yang2024leandojo} demonstrating the ability to generate proof steps in Lean. Our work is complementary: while these systems focus on proving mathematical theorems, Sparkle leverages Lean's proof capabilities to verify hardware properties. The LLM agent in our framework can potentially generate not only hardware implementations but also formal correctness proofs, a direction we explore through architecture equivalence verification.
